% Appendix: the model answers to Paper B, from exam/paper_b_solutions.md.

% The answers' headings are labelled ansb:q:N.
\renewcommand{\answerprefix}{ansb}

\chapter{Model Answers to Paper B}
\label{ansb:ch:answers}

Marks are shown per part. Method carries them: a correct design with a clumsy fix is worth more than
a correct fix with no reasoning, and several later parts consume earlier answers, so an error should
be followed through rather than penalised twice.

Code is marked on semantics. Missing includes, missing semicolons and misplaced braces cost nothing.

% ==========================================================================================
\answersection{1}{The workspace and the build}{12}

\answerpart{a}{4}

\begin{yamlcode}
manifest:
  remotes:
    - name: zephyrproject-rtos
      url-base: https://github.com/zephyrproject-rtos

  projects:
    - name: zephyr
      remote: zephyrproject-rtos
      revision: v4.4.2
      path: deps/zephyr
      clone-depth: 1
      west-commands: scripts/west-commands.yml

  self:
    path: .
\end{yamlcode}

\noindent\worth{2}

\noindent\emph{One mark for a manifest with a remote and a pinned project, one for \code{path},
\code{clone-depth} and \code{west-commands} together. \code{self: path: .} earns its mark below.}

Credit each of these where the candidate says it:
\begin{itemize}
  \item \textbf{\code{revision}} may be a tag, a branch or a commit hash; a product pins a tag or a
    hash, so that the same manifest checks out the same code next year.
  \item \textbf{\code{clone-depth: 1}} fetches only the pinned commit, which is most of the
    download; the cost is that \code{git log} in \file{deps/zephyr} shows one commit.
  \item \textbf{\code{west-commands}} is what makes \code{west build} and \code{west twister} exist.
    Without it west has only its own commands, because the book's manifest does not import Zephyr's.
  \item \textbf{\code{self: path: .}} says the manifest repository \emph{is} the workspace
    directory.
\end{itemize}

\textbf{What \code{west update} does:} it checks out every project in the manifest at the revision
the manifest names, cloning it first if it is not there. \worth{1}

\textbf{What it does not do:} it does not update the manifest. Nothing fetches a newer Zephyr by
itself: somebody edits \code{revision}, or a \code{git pull} of the companion repository brings a
new manifest in, and only then does \code{west update} make the workspace match it. Accept also:
it does not build anything, and it does not touch projects already at the right revision. \worth{1}

\answerpart{b}{4}

\begin{cmakecode}
cmake_minimum_required(VERSION 3.20.0)
find_package(Zephyr REQUIRED HINTS $ENV{ZEPHYR_BASE})
project(ready)

target_sources(app PRIVATE src/main.c)
\end{cmakecode}

\begin{kconfigcode}
CONFIG_LOG=y
\end{kconfigcode}

\begin{ccode}
#include <zephyr/kernel.h>
#include <zephyr/logging/log.h>

LOG_MODULE_REGISTER(app, LOG_LEVEL_INF);

int main(void)
{
	LOG_INF("ready");
	return 0;
}
\end{ccode}

\noindent\worth{2}

\noindent\emph{One mark for the three CMake lines plus \code{target_sources}, one for the module
registration and the log call. \file{prj.conf} must exist, and \code{CONFIG_LOG=y} must be in it;
an application with no \file{prj.conf} fails to configure with}
\code{No prj.conf file(s) was found}.

\textbf{What \code{find_package(Zephyr)} does.} It is where Zephyr's build system takes over: it
finds Zephyr through \code{ZEPHYR_BASE}, reads the board, processes the devicetree, runs Kconfig,
sets Zephyr up as a set of CMake libraries, and creates the library target \code{app} that the
application's own sources are added to. \worth{1}

\textbf{Why \code{project()} comes after it.} Zephyr selects the toolchain during
\code{find_package}, and CMake's \code{project()} is what makes a compiler current. Calling
\code{project()} first would have CMake look for a host compiler before Zephyr had said which
compiler to use. \worth{1}

\answerpart{c}{4}

\begin{shell}
west build -p -b native_sim/native/64 app
west build -t run
\end{shell}

\noindent\worth{1}

\noindent\emph{\code{-p} is pristine: it deletes the build directory first. Accept}
\code{./build/zephyr/zephyr.exe} \emph{in place of} \code{west build -t run}.

\begin{shell}
west build -p -b native_sim/native/64 app -- -DEXTRA_CONF_FILE=debug.conf
\end{shell}

\noindent\worth{1}

\noindent\emph{Everything after \code{--} goes to CMake. \code{CONF_FILE} in place of
\code{EXTRA_CONF_FILE} is wrong and earns nothing: it replaces \file{prj.conf} rather than merging
on top of it.}

\textbf{Three files}, any three of:

\begin{booktable}{@{}LL@{}}
  \thead{File} & \thead{What it holds} \\
  \midrule
  \file{build/zephyr/.config} & Every Kconfig option and the value the build settled on. \\
  \file{build/zephyr/include/generated/zephyr/autoconf.h} & The same configuration as C macros,
    included into every source file. \\
  \file{build/zephyr/zephyr.dts} & The final devicetree, overlays applied. \\
  \file{build/zephyr/include/generated/zephyr/devicetree_generated.h} & The devicetree as C
    macros. \\
  \file{build/zephyr/zephyr.exe} & The program itself, on \code{native_sim}. \\
  \file{build/zephyr/zephyr.map} & The linker map: where every symbol ended up. \\
  \file{build/compile_commands.json} & Every compiler command line, for editors. \\
\end{booktable}

\noindent\worth{2}

% ==========================================================================================
\answersection{2}{Configuring an application}{12}

\answerpart{a}{5}

\begin{kconfigcode}
menu "Application"

config APP_SAMPLE_PERIOD_MS
	int "Milliseconds between samples"
	range 10 10000
	default 1000

config APP_ALARM
	bool "Warn when a reading reaches the limit"

config APP_ALARM_LIMIT_C
	int "Temperature at which the alarm warns, in degrees Celsius"
	depends on APP_ALARM
	default 30

endmenu

source "Kconfig.zephyr"
\end{kconfigcode}

\noindent\worth{3}

\noindent\emph{One mark for the two typed options with prompt, range and defaults; one for
\code{depends on APP_ALARM} on the limit; one for \code{source "Kconfig.zephyr"} at the end. An
application's \file{Kconfig} replaces Zephyr's top-level one, so without that last line not one of
Zephyr's own options exists and the build fails on the board's own defaults, with}
\code{attempt to assign the value 'y' to the undefined symbol CONSOLE}.

\begin{kconfigcode}
CONFIG_APP_SAMPLE_PERIOD_MS=250
CONFIG_APP_ALARM=y
\end{kconfigcode}

\noindent\worth{1}

\textbf{With the alarm off, setting the limit aborts the build.} \code{APP_ALARM_LIMIT_C} depends on
\code{APP_ALARM}, so with the alarm off the symbol does not exist and the assignment cannot take.
Kconfig says so, and Zephyr treats the warning as an error:

\begin{console}
warning: APP_ALARM_LIMIT_C (defined at Kconfig:11) was assigned the value '35' but got the
value ''. Check these unsatisfied dependencies: APP_ALARM (=n) ...

error: Aborting due to Kconfig warnings
\end{console}

\noindent The alternative, a build that silently ignores the line, is what the abort exists to
prevent. \worth{1}

\answerpart{b}{4}

\begin{cmakecode}
target_sources_ifdef(CONFIG_APP_ALARM app PRIVATE src/alarm.c)
\end{cmakecode}

\noindent\worth{1}

\noindent\emph{Kconfig has run by the time \code{find_package(Zephyr)} returns, so every option is a
CMake variable too; this is how nearly all of Zephyr is built.}

\begin{ccode}
LOG_MODULE_REGISTER(sampler, LOG_LEVEL_INF);

/* ... */
	LOG_INF("Reading %d C", celsius);

	if (IS_ENABLED(CONFIG_APP_ALARM)) {
		alarm_check(celsius);
	}
\end{ccode}

\noindent with \file{alarm.c} holding

\begin{ccode}
LOG_MODULE_DECLARE(sampler);

void alarm_check(int celsius)
{
	if (celsius >= CONFIG_APP_ALARM_LIMIT_C) {
		LOG_WRN("Reading %d C is at or above the limit of %d C", celsius,
			CONFIG_APP_ALARM_LIMIT_C);
	}
}
\end{ccode}

\noindent\worth{2}

\noindent\emph{One mark for the module registration and the guarded call, one for
\code{LOG_MODULE_DECLARE} in the second file. A second \code{LOG_MODULE_REGISTER} of the same name
is the mistake here; the module is registered once and declared everywhere else.}

\textbf{Why \code{IS_ENABLED()}.} It expands to \code{1} or \code{0}, so the call is compiled and
type-checked in every configuration, and the compiler removes the branch that cannot run --- call
included, which is why the program still links when \file{alarm.c} is not compiled. Code behind
\code{#ifdef} is not even parsed when the option is off, so it rots unnoticed until somebody builds
the other configuration. \worth{1}

\answerpart{c}{3}
\textbf{In order:}
\begin{enumerate}
  \item the option's own \code{default} in the \file{Kconfig} file;
  \item the board's defaults, from its \code{_defconfig};
  \item the application's \file{prj.conf}, and any board-specific \file{.conf} beside it;
  \item extra configuration files given with \code{-DEXTRA_CONF_FILE=};
  \item \code{-DCONFIG_...} on the \code{west build} command line.
\end{enumerate}

\noindent\textbf{The last one applied wins}, so the command line beats everything else. Accept any
four of these five in the right order. \worth{2}

\textbf{The second build still uses 50 ms.} A \code{-DCONFIG_...} option is stored in the build
directory's CMake cache, and every later build in that directory keeps using it until the directory
is rebuilt with \code{-p} or the option is given a new value. This is one of the ways a build
directory goes stale, and why a pristine build is the only way to be sure that nothing but the files
and the command line decide the configuration. \worth{1}

% ==========================================================================================
\answersection{3}{Describing hardware}{13}

\answerpart{a}{5}

\begin{yamlcode}
description: |
  A heater driven by one GPIO pin.

compatible: "qacademy,heater"

properties:
  control-gpios:
    type: phandle-array
    required: true
    description: The pin that switches the heater on.

  zone:
    type: string
    required: true
    description: The name of the zone this heater warms.

  max-power-percent:
    type: int
    default: 100
    description: The highest duty cycle this heater may be driven at.
\end{yamlcode}

\noindent\worth{3}

\noindent\emph{One mark for \code{compatible} and the three property names, one for the types, one
for \code{required: true} on two and \code{default: 100} on the third. A \code{default} on a
required property is a contradiction and loses the last mark.}

\textbf{Where it goes:} \file{dts/bindings/} in the application directory --- the application's own
directory is searched for bindings, alongside Zephyr's \file{dts/bindings}. The file name does not
matter to the build, but the convention is the compatible:
\file{dts/bindings/qacademy,heater.yaml}. \worth{1}

\textbf{The vendor prefix.} \code{qacademy} is not one of the prefixes Zephyr knows, so the build
warns \code{unknown vendor prefix 'qacademy'}. A \file{dts/bindings/vendor-prefixes.txt} beside the
binding, with the line \code{qacademy} and a tab and a description, registers it. This matters more
than the warning suggests: Twister builds with \code{--edtlib-Werror}, where the same warning is an
error. \worth{1}

\answerpart{b}{4}

\begin{dtscode}
/ {
	heater_a: heater-a {
		compatible = "qacademy,heater";
		control-gpios = <&gpio0 2 GPIO_ACTIVE_HIGH>;
		zone = "cabin";
	};

	heater_b: heater-b {
		compatible = "qacademy,heater";
		control-gpios = <&gpio0 3 GPIO_ACTIVE_LOW>;
		zone = "tank";
		max-power-percent = <60>;
	};

	aliases {
		primary-heater = &heater_a;
	};
};
\end{dtscode}

\noindent\worth{3}

\noindent\emph{One mark for two nodes with the right compatible at the root, one for the
\code{gpios} cells with their flags, one for the alias. \code{max-power-percent} must be left out of
\code{heater_a}, which is what makes the binding's default do its work.}

\textbf{The file name is \file{app.overlay}}, in the application directory: the build picks that up
with no argument and no CMake line. Accept \file{boards/native_sim_native_64.overlay} with the note
that a board-named overlay \emph{replaces} \file{app.overlay} rather than adding to it.

\textbf{What it does:} the overlay is merged on top of the board's devicetree before anything is
generated from it, adding nodes, changing properties, or disabling nodes the board declared. The
board's own file is not edited; the result is what \file{build/zephyr/zephyr.dts} shows. \worth{1}

\answerpart{c}{4}

\begin{ccode}
struct heater {
	struct gpio_dt_spec control;
	const char *zone;
	int max_power_percent;
};

#define HEATER_ENTRY(node)                                                        \
	{                                                                         \
		.control = GPIO_DT_SPEC_GET(node, control_gpios),                 \
		.zone = DT_PROP(node, zone),                                      \
		.max_power_percent = DT_PROP(node, max_power_percent),            \
	},

static const struct heater heaters[] = {DT_FOREACH_STATUS_OKAY(qacademy_heater, HEATER_ENTRY)};

int main(void)
{
	for (size_t i = 0; i < ARRAY_SIZE(heaters); i++) {
		const struct heater *h = &heaters[i];

		if (!gpio_is_ready_dt(&h->control)) {
			LOG_ERR("%s: controller not ready", h->zone);
			continue;
		}
		gpio_pin_configure_dt(&h->control, GPIO_OUTPUT_INACTIVE);
		LOG_INF("%-6s pin %u, at most %d %%", h->zone, h->control.pin,
			h->max_power_percent);
	}
	return 0;
}
\end{ccode}

\noindent\worth{3}

\noindent\emph{One mark for \code{DT_FOREACH_STATUS_OKAY} with the compatible spelled in C, commas
and dashes replaced by underscores, and a macro that expands to one initializer including its
trailing comma; one for \code{GPIO_DT_SPEC_GET} and \code{DT_PROP}; one for the readiness check and
configuring the pin as an inactive output. The table is \code{static const}: it is built by the
compiler and lives in flash. Nothing here is a search at run time.}

\textbf{With \code{status = "disabled"}} the node is not ``status okay'', so
\code{DT_FOREACH_STATUS_OKAY} does not expand it: the table has one entry, \code{ARRAY_SIZE} is 1,
and nothing else changes. It still builds.

\textbf{With \code{zone} deleted} the build stops during configuration, before any C is compiled:

\begin{console}
devicetree error: 'zone' is marked as required in 'properties:' in
'.../dts/bindings/qacademy,heater.yaml', but does not appear in <Node /heater-b in
'.../app.overlay'>
\end{console}

\noindent That is the binding doing its job: a required property is a promise the devicetree must
keep. \worth{1}

% ==========================================================================================
\answersection{4}{Threads and time}{13}

\answerpart{a}{5}
A defensible answer, and the one the book leads to:

\newcolumntype{R}[1]{>{\raggedright\arraybackslash}p{#1}}
\begin{booktable}{@{}R{7em}R{3.4em}R{5.4em}L@{}}
  \thead{Activity} & \thead{Priority} & \thead{Kind} & \thead{Why} \\
  \midrule
  Control loop & \code{-1} & Cooperative & It must not be late. A cooperative thread cannot be
    preempted at all, so once it is running nothing else runs until it blocks --- which it does, on
    its next period. \\
  Display & \code{5} & Preemptible & It must not delay the control loop, and it does not mind being
    interrupted part-way through a refresh. \\
  Logger & \code{7} & Preemptible & It is the one that may be late. At the lowest priority it runs
    only when the other two are waiting, and its 50 ms of flash work delays nothing. \\
\end{booktable}

\begin{ccode}
K_THREAD_DEFINE(control, 1024, control_thread, NULL, NULL, NULL, -1, 0, 0);
K_THREAD_DEFINE(display, 1024, display_thread, NULL, NULL, NULL, 5, 0, 0);
K_THREAD_DEFINE(logger, 2048, logger_thread, NULL, NULL, NULL, 7, 0, 0);
\end{ccode}

\noindent\worth{5}

\noindent\emph{Three marks for the three choices with a reason each; two for the three definitions.
In Zephyr a \textbf{lower number is a higher priority}, and negative priorities are the cooperative
ones: getting that backwards costs the marks for the choices. Any consistent set of preemptible
priorities is fine --- what is marked is control above display above logger, and the reasoning.}

Credit as an addition: the logger is the thread to give the larger stack, because it is the one that
calls into a file system or flash; and a cooperative control loop is only safe because its body is
short and ends in a wait.

\answerpart{b}{4}

\begin{ccode}
static K_TIMER_DEFINE(control_tick, NULL, NULL);

static void control_thread(void)
{
	k_timer_start(&control_tick, K_MSEC(10), K_MSEC(10));

	while (true) {
		k_timer_status_sync(&control_tick);   /* Waits for the next expiry. */
		control_step();
	}
}
\end{ccode}

\noindent\worth{2}

\noindent\emph{\code{k_timer_status_sync()} blocks until the timer has expired at least once and
returns how many times it has, so a slow iteration shows up as a count above 1 rather than as
drift.}

\textbf{Why \code{k_msleep(10)} drifts.} It sleeps for 10 ms \emph{from the moment it is called},
which is after the body has run. Every iteration is therefore the body plus 10 ms plus the rounding
to the next tick, and the error accumulates: the loop's period is never 10 ms and the tenth
iteration is late by ten roundings. A periodic timer measures from the previous expiry, so its
errors do not add up. \worth{1}

\textbf{A second way:} absolute timeouts. Keep a deadline and sleep until it, rather than for a
duration:

\begin{ccode}
int64_t next = k_uptime_get() + 10;

while (true) {
	k_sleep(K_TIMEOUT_ABS_MS(next));
	control_step();
	next += 10;
}
\end{ccode}

\noindent This is equivalent because the wake-up time, not the sleep's length, is what is specified:
the body's duration comes out of the wait rather than being added to it. \worth{1}

\answerpart{c}{4}
\textbf{\code{k_msleep(30)} returns after 40 ms.} The kernel keeps timeouts in ticks, and
\code{native_sim} has 100 of them a second, so 30 ms is three ticks. The thread calls at some point
\emph{within} the current tick, and the kernel cannot know how much of it is left, so it waits the
three ticks the duration needs plus the remainder of the one in progress: four tick boundaries,
40 ms. A sleep is a \emph{minimum}. \worth{2}

\textbf{The timer expires at 30, 60 and 90.} \code{k_timer_start()} compensates for exactly that
added tick, and a periodic timer's next expiry is measured from the previous one rather than from
when the handler ran, so a periodic timer does not drift. \worth{1}

\textbf{With 1000 ticks per second}, a tick is 1 ms, so the same sleep returns after 31 ms: the
rounding is still there but is ten times smaller. The cost is that the kernel's clock bookkeeping
happens ten times as often; on a board that is more timer interrupts, more power, and less time
asleep. \worth{1}

% ==========================================================================================
\answersection{5}{Interrupts, work and the simulator}{12}

\answerpart{a}{5}

\begin{ccode}
static void sample_work_handler(struct k_work *work);

static K_WORK_DEFINE(sample_work, sample_work_handler);

static void sample_timer_expired(struct k_timer *timer)
{
	ARG_UNUSED(timer);
	k_work_submit(&sample_work);          /* All an ISR should do: hand the work on. */
}

static K_TIMER_DEFINE(sample_timer, sample_timer_expired, NULL);

static void sample_work_handler(struct k_work *work)
{
	ARG_UNUSED(work);
	LOG_INF("sample %d", read_sensor());
}

int main(void)
{
	k_timer_start(&sample_timer, K_MSEC(CONFIG_APP_SAMPLE_PERIOD_MS),
		      K_MSEC(CONFIG_APP_SAMPLE_PERIOD_MS));
	return 0;
}
\end{ccode}

\noindent\worth{3}

\noindent\emph{One mark for the timer with a period, one for the expiry function doing nothing but
submitting, one for the work item and its handler.}

\textbf{Two things the expiry function must not do:}
\begin{itemize}
  \item \textbf{Sleep or block} --- \code{k_sleep()}, \code{k_mutex_lock()} with a timeout,
    \code{k_msgq_get(..., K_FOREVER)}. A timer's expiry function runs in the system clock's
    interrupt handler, and there is no thread to put to sleep. With \code{CONFIG_ASSERT=y} the
    kernel catches it and panics; without, the system misbehaves.
  \item \textbf{Take long} --- the expiry function delays every interrupt behind it and every
    thread. Anything that reads a sensor over a bus, writes to flash, or formats a long log line
    belongs in the work handler.
\end{itemize}

\noindent\worth{1}

\textbf{Which thread the handler runs in:} the system work queue's thread, \code{sysworkq}, a
cooperative thread shared by every subsystem. That is why it may sleep and log, and why it should
still be quick: everybody else's work is queued behind it. A queue of your own, started with
\code{k_work_queue_start()}, is the answer when the work is slow. \worth{1}

\answerpart{b}{4}

\begin{ccode}
static void idle_work_handler(struct k_work *work)
{
	ARG_UNUSED(work);
	LOG_INF("idle");
}

static K_WORK_DELAYABLE_DEFINE(idle_work, idle_work_handler);
\end{ccode}

\noindent and, at the end of the sample handler:

\begin{ccode}
	k_work_reschedule(&idle_work, K_MSEC(250));
\end{ccode}

\noindent\worth{2}

\textbf{Why \code{k_work_reschedule()}.} It \emph{moves} an already scheduled item to the new time;
\code{k_work_schedule()} leaves an already scheduled item alone. With \code{schedule}, the timeout
would keep the time set by the first sample and fire 250 ms after \emph{that}, whatever happened
since. Pushing the deadline back on every sample is exactly what ``no sample for 250 ms'' means.
\worth{1}

\textbf{When it fires.} 250 ms is 25 ticks, plus the remainder of the tick in progress: 26 ticks, so
260 ms after the last sample. With the last sample at 500 ms, the message is logged at
\textbf{760 ms}. Accept ``at 750 ms, give or take a tick'' with the rounding named. \worth{1}

\answerpart{c}{3}
\textbf{It hangs, and takes the whole system with it.} On \code{native_sim} the code runs in zero
simulated time and interrupts arrive only when the CPU is idle, which it never becomes while a
thread spins. The timer's expiry function is never called, \code{ready} is never set, the clock
never advances, and even \code{-stop_at} cannot end the run, because that too is a simulated-time
event. On a real microcontroller the same loop would end --- the timer interrupt would preempt it
--- but it would burn the CPU for the whole wait. \worth{2}

\textbf{The fix} is to wait on something rather than to look at it:

\begin{ccode}
static K_SEM_DEFINE(ready, 0, 1);

static void sample_timer_expired(struct k_timer *timer)
{
	k_sem_give(&ready);
}

/* In the thread: */
	k_sem_take(&ready, K_FOREVER);
\end{ccode}

\noindent\worth{1}

\noindent\emph{Accept \code{k_poll}, \code{k_timer_status_sync()}, or a \code{k_msgq} or zbus wait.
A \code{k_msleep()} inside the loop is worth half: it lets time pass, so it works, but it is
polling.}

% ==========================================================================================
\answersection{6}{Threads that must agree}{13}

\answerpart{a}{5}

\begin{ccode}
struct sample {
	uint32_t seq;
	int32_t celsius;
};

K_MSGQ_DEFINE(samples, sizeof(struct sample), 4, alignof(struct sample));

/* The sampler, every 100 ms: */
	struct sample s = {.seq = ++seq, .celsius = read_sensor()};

	if (k_msgq_put(&samples, &s, K_NO_WAIT) != 0) {
		dropped++;
	}

/* The storage thread: */
	struct sample s;

	while (k_msgq_get(&samples, &s, K_FOREVER) == 0) {
		store(&s);
	}
\end{ccode}

\noindent\worth{4}

\noindent\emph{One mark for \code{K_MSGQ_DEFINE} with the message size, the depth and the alignment;
one for \code{K_NO_WAIT} and counting the failure; one for \code{K_FOREVER} in the consumer; one for
checking both return values. A producer that passes a timeout, or ignores the result of
\code{k_msgq_put()}, has silently dropped the sample instead of counting it.}

\textbf{What \code{k_msgq_put()} does with the struct:} it copies the bytes into the queue's own
storage, which was allocated when the queue was defined. Nothing is shared and nothing is owned, so
the producer's local may go out of scope the moment the call returns --- and a reader gets a
snapshot, not a reference to something that may have changed since. \worth{1}

\answerpart{b}{4}

\begin{ccode}
static K_MUTEX_DEFINE(calibration_lock);

	k_mutex_lock(&calibration_lock, K_FOREVER);
	/* ... the critical section ... */
	k_mutex_unlock(&calibration_lock);
\end{ccode}

\noindent\worth{1}

\textbf{Two things a mutex gives you that a binary semaphore does not:}
\begin{itemize}
  \item \textbf{Ownership.} A mutex is held by a thread, and only that thread can unlock it; it is
    also recursive for its owner. A semaphore has no owner: any thread can give it, including one
    that never took it, so an unbalanced give goes unnoticed.
  \item \textbf{Priority inheritance.} While a higher-priority thread waits for the mutex, the
    kernel raises the holder's priority to the waiter's, and restores it at the unlock.
\end{itemize}

\noindent Credit also: a mutex may not be used in an ISR, which is a useful restriction, since an
ISR has no priority to inherit and no identity to own with. \worth{2}

\textbf{What the semaphore version suffers from: unbounded priority inversion.} The thread at 7
holds the guard; the one at 5 blocks on it; the one at 6, which never touches the guard, preempts 7
and runs for as long as it likes. The highest-priority thread in the system waits for the middle
one, and its wait is bounded by nothing anybody designed.

\textbf{With the mutex}, 7 runs at 5's priority while it holds the lock, so 6 cannot preempt it. The
highest-priority thread's wait is bounded by the length of the critical section --- which is the
thing the designer controls, and the reason to keep it short. \worth{1}

\answerpart{c}{4}

\begin{ccode}
static bool sample_is_valid(const void *msg, size_t msg_size)
{
	const struct sample *s = msg;

	ARG_UNUSED(msg_size);
	return (s->celsius >= -40) && (s->celsius <= 125);
}

static void alarm_callback(const struct zbus_channel *chan)
{
	const struct sample *s = zbus_chan_const_msg(chan);

	LOG_INF("alarm saw #%u", s->seq);
}

ZBUS_LISTENER_DEFINE(alarm, alarm_callback);
ZBUS_SUBSCRIBER_DEFINE(display, 4);

ZBUS_CHAN_DEFINE(sample_chan, struct sample, sample_is_valid, NULL,
		 ZBUS_OBSERVERS(alarm, display), ZBUS_MSG_INIT(0));

static void display_thread(void)
{
	const struct zbus_channel *chan;
	struct sample s;

	while (zbus_sub_wait(&display, &chan, K_FOREVER) == 0) {
		zbus_chan_read(chan, &s, K_FOREVER);
		LOG_INF("display saw #%u", s.seq);
	}
}
\end{ccode}

\noindent\worth{3}

\noindent\emph{One mark for the channel with its message type and validator, one for the listener,
one for the subscriber and its wait-then-read loop. The publisher calls
\code{zbus_chan_pub(&sample_chan, &s, ...)}; a rejected message gives it \code{-ENOMSG}.}

\textbf{Where each one runs.} The \textbf{listener} is a callback run inside
\code{zbus_chan_pub()}, in the publisher's own thread and with the channel locked --- so it must be
quick and must not block, like an ISR. The \textbf{subscriber} runs in its own thread, whenever the
scheduler gets to it.

\textbf{Why a subscriber can miss a message.} It is not handed the message: it is handed a
notification that the channel changed, and reads the channel later, by which time the channel holds
whatever was published last. Two publications before the subscriber runs leave it reading the second
one twice, or seeing only the newer value. The listener cannot miss one, because it is called during
the publication, before \code{zbus_chan_pub()} returns. Credit the message subscriber,
\code{ZBUS_MSG_SUBSCRIBER_DEFINE}, as the observer that keeps a copy of every message and so misses
none. \worth{1}

% ==========================================================================================
\answersection{7}{A driver of your own}{13}

\answerpart{a}{6}

\begin{ccode}
#define DT_DRV_COMPAT qacademy_baro

#include <zephyr/device.h>
#include <zephyr/drivers/i2c.h>
#include <zephyr/drivers/sensor.h>
#include <zephyr/logging/log.h>

#include "baro_reg.h"

LOG_MODULE_REGISTER(baro, CONFIG_SENSOR_LOG_LEVEL);

struct baro_config {
	struct i2c_dt_spec i2c;
};

struct baro_data {
	uint16_t raw;
};

static int baro_sample_fetch(const struct device *dev, enum sensor_channel chan)
{
	const struct baro_config *config = dev->config;
	struct baro_data *data = dev->data;
	uint8_t buf[2];
	int ret;

	if ((chan != SENSOR_CHAN_ALL) && (chan != SENSOR_CHAN_PRESS)) {
		return -ENOTSUP;
	}

	ret = i2c_burst_read_dt(&config->i2c, BARO_REG_PRESSURE_L, buf, sizeof(buf));
	if (ret < 0) {
		return ret;
	}

	data->raw = (uint16_t)buf[0] | ((uint16_t)buf[1] << 8);
	return 0;
}

static int baro_channel_get(const struct device *dev, enum sensor_channel chan,
			    struct sensor_value *val)
{
	const struct baro_data *data = dev->data;

	if (chan != SENSOR_CHAN_PRESS) {
		return -ENOTSUP;
	}

	/* One count is 10 Pa; the sensor API's pressure channel is in kilopascals. */
	val->val1 = data->raw / 100;
	val->val2 = (data->raw % 100) * 10000;
	return 0;
}

static DEVICE_API(sensor, baro_api) = {
	.sample_fetch = baro_sample_fetch,
	.channel_get = baro_channel_get,
};

static int baro_init(const struct device *dev)
{
	const struct baro_config *config = dev->config;
	uint8_t id;
	int ret;

	if (!i2c_is_ready_dt(&config->i2c)) {
		return -ENODEV;
	}

	ret = i2c_reg_read_byte_dt(&config->i2c, BARO_REG_WHO_AM_I, &id);
	if (ret < 0) {
		return ret;
	}

	if (id != BARO_WHO_AM_I_VALUE) {
		LOG_ERR("%s: WHO_AM_I is 0x%02x, expected 0x%02x", dev->name, id,
			BARO_WHO_AM_I_VALUE);
		return -ENODEV;
	}

	return 0;
}

#define BARO_DEFINE(inst)                                                         \
	static struct baro_data baro_data_##inst;                                 \
	static const struct baro_config baro_config_##inst = {                    \
		.i2c = I2C_DT_SPEC_INST_GET(inst),                                \
	};                                                                        \
	SENSOR_DEVICE_DT_INST_DEFINE(inst, baro_init, NULL, &baro_data_##inst,    \
				     &baro_config_##inst, POST_KERNEL,            \
				     CONFIG_SENSOR_INIT_PRIORITY, &baro_api);

DT_INST_FOREACH_STATUS_OKAY(BARO_DEFINE)
\end{ccode}

\noindent\worth{6}

One mark each for:
\begin{enumerate}
  \item \textbf{\code{DT_DRV_COMPAT}} spelled with the comma turned into an underscore, before the
    includes. Every \code{DT_INST_} macro in the file counts instances of that compatible only.
  \item \textbf{The config/data split}: what the devicetree says is \code{const} and lives in
    \code{dev->config}; what changes at run time lives in \code{dev->data}. A driver that keeps its
    \code{i2c_dt_spec} in \code{data}, or that makes the config non-const, loses this mark --- on a
    board the config is in flash.
  \item \textbf{\code{baro_init} checking the bus and the identity}, and returning \code{-ENODEV}
    rather than 0 when the part answers with the wrong value.
  \item \textbf{\code{sample_fetch}} reading both bytes in one burst and storing the raw value, with
    the channel checked.
  \item \textbf{\code{channel_get}} converting to \code{sensor_value} correctly: 10132 counts gives
    \code{val1} 101 and \code{val2} 320000, because \code{val2} is millionths.
  \item \textbf{\code{DT_INST_FOREACH_STATUS_OKAY(BARO_DEFINE)}} with
    \code{I2C_DT_SPEC_INST_GET}, so that the driver serves every node with that compatible rather
    than one.
\end{enumerate}

\noindent\emph{A note for the marker: \code{SENSOR_DEVICE_DT_INST_DEFINE} and
\code{DEVICE_DT_INST_DEFINE} are both acceptable, as is a plain
\code{static const struct sensor_driver_api} in place of \code{DEVICE_API}.}

\answerpart{b}{4}

\begin{yamlcode}
name: qacademy-baro
build:
  cmake: .
  kconfig: Kconfig
  settings:
    dts_root: .
\end{yamlcode}

\noindent\worth{1}

\noindent\emph{\code{dts_root: .} is what makes the module's \file{dts/bindings} visible to the
build; without it the binding is not found and every node with that compatible has no binding.}

\begin{kconfigcode}
config BARO
	bool "Baro I2C pressure sensor"
	default y
	depends on SENSOR
	depends on DT_HAS_QACADEMY_BARO_ENABLED
	select I2C
\end{kconfigcode}

\noindent\worth{1}

\begin{yamlcode}
description: Baro, a made-up I2C pressure sensor reading in steps of 10 Pa.

compatible: "qacademy,baro"

include: [sensor-device.yaml, i2c-device.yaml]
\end{yamlcode}

\noindent\worth{1}

\noindent\emph{\file{i2c-device.yaml} is what requires \code{reg}, the bus address, and makes the
node an I2C device; \file{sensor-device.yaml} brings the properties every sensor may have.}

\textbf{Why the dependency.} \code{DT_HAS_QACADEMY_BARO_ENABLED} is generated from the devicetree:
it is defined exactly when some node has that compatible and status okay. Tying the driver to it
means the driver is built when, and only when, the hardware is described --- which is why enabling a
driver in Zephyr usually means adding a devicetree node rather than setting an option. An
application without the part gets none of the driver's code, and \code{default y} costs it nothing.
\worth{1}

\answerpart{c}{3}

\begin{ccode}
#define BARO_NODE DT_NODELABEL(baro)

static const struct device *const baro = DEVICE_DT_GET(BARO_NODE);
static const struct emul *const baro_emul = EMUL_DT_GET(BARO_NODE);

ZTEST_SUITE(baro, NULL, NULL, NULL, NULL, NULL);

ZTEST(baro, test_device_is_ready)
{
	zassert_true(device_is_ready(baro));
}

ZTEST(baro, test_reading_converts_to_kilopascals)
{
	struct sensor_value pressure;

	baro_emul_set_raw(baro_emul, 10132);
	zassert_ok(sensor_sample_fetch(baro));
	zassert_ok(sensor_channel_get(baro, SENSOR_CHAN_PRESS, &pressure));
	zassert_equal(pressure.val1, 101);
	zassert_equal(pressure.val2, 320000);
}
\end{ccode}

\noindent\worth{2}

\noindent\emph{One mark for the suite and the readiness test, one for driving the emulator and
asserting on both halves of the \code{sensor_value}.}

\textbf{The file is \file{testcase.yaml}}, beside the test application, naming the platform and
\code{harness: ztest}. \code{west twister -T <dir> -p native_sim/native/64} then builds and runs it.

\textbf{\code{device_is_ready()} returns false} for a device whose init returned a negative error:
the kernel records the failure and the device stays unusable for the life of the program. The
application must therefore check it before the first call --- that is what the check is for, and a
driver that fails its identity check is exactly the case it catches. \worth{1}

% ==========================================================================================
\answersection{8}{Making it usable}{12}

\answerpart{a}{5}

\begin{ccode}
static int32_t setpoint_c = 20;

static int cmd_setpoint_get(const struct shell *sh, size_t argc, char **argv)
{
	ARG_UNUSED(argc);
	ARG_UNUSED(argv);

	shell_print(sh, "setpoint %d C", setpoint_c);
	return 0;
}

static int cmd_setpoint_set(const struct shell *sh, size_t argc, char **argv)
{
	char *end;
	long value = strtol(argv[1], &end, 10);

	ARG_UNUSED(argc);

	if ((*end != '\0') || (value < -40) || (value > 125)) {
		shell_error(sh, "setpoint must be -40..125 degrees Celsius, not '%s'", argv[1]);
		return -EINVAL;
	}

	setpoint_c = (int32_t)value;
	shell_print(sh, "setpoint %d C", setpoint_c);
	return 0;
}

SHELL_STATIC_SUBCMD_SET_CREATE(setpoint_cmds,
			       SHELL_CMD(get, NULL, "Print the setpoint.", cmd_setpoint_get),
			       SHELL_CMD_ARG(set, NULL, "Set the setpoint: set <celsius>.",
					     cmd_setpoint_set, 2, 0),
			       SHELL_SUBCMD_SET_END);

SHELL_CMD_REGISTER(setpoint, &setpoint_cmds, "The temperature setpoint.", NULL);
\end{ccode}

\noindent\worth{4}

\noindent\emph{One mark for the handler signature and \code{shell_print} to the shell that ran the
command rather than \code{printk}; one for the subcommand set and its \code{SHELL_SUBCMD_SET_END};
one for \code{SHELL_CMD_ARG} with two mandatory words; one for validating the text as well as the
range --- \code{strtol} alone accepts \code{abc} as 0, which is why the end pointer is checked.}

\textbf{What the shell has already checked:} the argument \emph{count}.
\code{SHELL_CMD_ARG(..., 2, 0)} declares two mandatory words and no optional ones, so
\code{setpoint set} on its own never reaches the handler; the shell rejects it with
\code{set: wrong parameter count} and prints the command's help. What is left to the handler is what
the words \emph{say}.

\textbf{The return value} is 0 for success or a negative errno; the shell reports a failure to the
user, and it is what a script or a test driving the shell sees. \worth{1}

\answerpart{b}{4}

\begin{ccode}
static int setpoint_set_cb(const char *name, size_t len, settings_read_cb read_cb, void *cb_arg)
{
	const char *next;

	if (settings_name_steq(name, "celsius", &next) && (next == NULL)) {
		if (len != sizeof(setpoint_c)) {
			return -EINVAL;
		}
		ssize_t read = read_cb(cb_arg, &setpoint_c, sizeof(setpoint_c));

		return (read < 0) ? (int)read : 0;
	}
	return -ENOENT;
}

SETTINGS_STATIC_HANDLER_DEFINE(setpoint, "setpoint", NULL, setpoint_set_cb, NULL, NULL);
\end{ccode}

\noindent with, at start-up, \code{settings_subsys_init()} then \code{settings_load()}, and in the
\code{set} handler:

\begin{ccode}
	settings_save_one("setpoint/celsius", &setpoint_c, sizeof(setpoint_c));
\end{ccode}

\noindent\worth{2}

\noindent\emph{One mark for a handler registered for the subtree and reading through
\code{read_cb} rather than expecting a pointer to the value, one for \code{settings_save_one()} with
the full name and \code{settings_subsys_init()}/\code{settings_load()} at boot.}

\textbf{The Kconfig options:} \code{CONFIG_SETTINGS=y}, a storage back-end ---
\code{CONFIG_ZMS=y} or \code{CONFIG_NVS=y} --- and the flash underneath it, \code{CONFIG_FLASH=y}
and \code{CONFIG_FLASH_MAP=y}. \worth{1}

\textbf{Where it is kept:} in the flash partition labelled \code{storage_partition} in the board's
devicetree. On \code{native_sim} that flash is a file, \file{flash.bin}, in the directory the
program is run from --- the build directory when it is started with \code{west build -t run}.

\textbf{Two ways to make the next run a first boot}, any two of: delete \file{flash.bin}; start the
program with \code{-flash_erase}; start it with \code{-flash_in_ram}, so the flash never touches the
disk; or build with \code{-p}, which deletes the build directory and the file with it. \worth{1}

\answerpart{c}{3}

\begin{ccode}
void k_sys_fatal_error_handler(unsigned int reason, const struct arch_esf *esf)
{
	ARG_UNUSED(esf);

	LOG_PANIC();
	LOG_ERR("fatal error %u in thread '%s'", reason, k_thread_name_get(k_current_get()));

	if (reason == K_ERR_KERNEL_OOPS) {
		return;                 /* The kernel aborts this thread only. */
	}

	k_fatal_halt(reason);
}
\end{ccode}

\noindent\worth{2}

\noindent\emph{One mark for replacing the weak handler with one that returns on
\code{K_ERR_KERNEL_OOPS}, one for halting otherwise. \code{LOG_PANIC()} earns credit as an addition:
it puts logging into synchronous mode, so that what has been logged is actually printed before the
system goes down --- which matters on a board, where logging is deferred to a thread that may never
run again.}

\textbf{The exit status is 1.} Halting on \code{native_sim} ends the process, and the harness sees a
failure; that is also why Twister treats such a run as failed whatever it printed.

\textbf{Why returning from a panic is not allowed:} an oops means one thread has found something
impossible, and the kernel can abort just that thread and carry on. A panic means the kernel itself
cannot continue --- a failed assertion inside it, a corrupted structure, a fault with no thread to
blame --- so there is nothing to return \emph{to}. The kernel's own handler halts for exactly that
reason. \worth{1}
